%% kapa: exam=1 round=23 year=2012 period=1 form=A subject=economics subject_ko=경제원론 seq=38 number=78
%% source: Q-Net 감정평가사 기출문제 (cst00309) artlSeq=5052578 file=raw/1008/5052578/2012년 제23회 감정평가사 제1차시험 기출문제(1교시 A형).pdf
%% conversion: kordoc 4.24.51 -> tools/gen1.py (auto); review=proofread-sonnet; flags=none
\begin{question}{78}
중앙은행의 통화량 조절 방법에 관한 설명으로 옳은 것은?
\begin{choices}
  \item 법정지급준비율을 인상하면 시중은행이 예금액 중에서 대출할 수 있는 금액이 증가한다.
  \item 중앙은행이 국채를 시중은행 A에 매도하면 시중은행 A의 지급준비금은 증가한다.
  \item 법정지급준비율을 인하하면 예금통화승수는 감소한다.
  \item 재할인율을 인상하면 통화량이 증가한다.
  \item 중앙은행이 민간인들이 보유하고 있는 국채를 매입하면 통화량은 증가한다.
\end{choices}
\end{question}
\answer{5}
